


This note records an exploratory comparison family of Gamma measures.
The independent order $K$ is a comparison parameter: the identities below
do not derive this order from the original tensor, exterior, or spectral-sum
operation. The original packet parameter $k$ and its canonical source
remain as specified by those original constructions.

\section{The independent order and the unchanged original coordinate}

Keep the original packet integers and centre
\[
 k=4l+1,\quad e=1+k(m-1),\quad q=e(k+1)^2=2n,
 \quad c=k/2,
\]
and choose a separate integer $L$ with $1\le L\le q/4$. Its Gamma
order is $K=4L+1$. The original packet order $k$ has not been replaced
by $K$, and the source coordinate remains $S=c+iy$. Since $q$ is
divisible by four, the exact choice $L=q/4$ gives $K=q+1$.
The original Gamma measure and the actual multiplying polynomial are
\[
 d\sigma(y)=\frac{|\Gamma(1/4+iy/2)|^2}{2\pi}\,dy,
 \quad \sigma(\R)=\sqrt{2\pi},\qquad
 f_L(S)=\prod_{j=0}^{L-1}(S-c-2j-1/2),
 \tag{MFG1}
\]
\[
 \Phi_L(y)=|f_L(c+iy)|^2
          =\prod_{j=0}^{L-1}\big(y^2+(2j+1/2)^2\big).
 \tag{MFG2}
\]
We first calculate the exact homogeneous relation source
$\chi_0(S)=(S-c)^q$. The original displaced packet returns by its
separate full-source comparison, with the same $f_L$ retained.

Put $\mu_a=\int y^{2a}d\sigma(y)$ and let $\nu=(y\mapsto y^2)_*\sigma$
be the full square pushforward, including both half-axes. For integers
$a,s\ge0$ define
\[
 H_s^{(a;L)}=\det\left[
    \int_0^\infty t^{a+i+j}
             \prod_{r=0}^{L-1}(t+(2r+1/2)^2)\,d\nu(t)
                  \right]_{i,j<s},\qquad H_0^{(a;L)}=1.
 \tag{MFG3}
\]
At $L=0$ these are exactly the old $H_s^{(a)}$. The full real-line
relation matrices have entries
\[\int y^{2q+i+j}\Phi_L(y)d\sigma(y).\] Their even/odd congruence gives
the literal paired determinant
\[
 Z_{q,L}=H_n^{(q;L)}H_{n+1}^{(q;L)}(H_n^{(q+1;L)})^2,
 \qquad Z_{q,0}=Z_q.
 \tag{MFG4}
\]
The same row and column permutation is used, so its determinant sign
is squared. The low relation ratio and the signed source product are
\[
 r_{q,L}=\frac{\int y^{2q}\Phi_L(y)d\sigma(y)}{\mu_q},
\]
\[
 P_{q,L}=-\sum_{j=0}^{2L-1}
 \left[\sum_{a=q}^{2q-1}\log(a+j+1/2)
       +\sum_{a=q+1}^{2q}\log(a+j+1/2)\right].
 \tag{MFG5}
\]
Here the complete source determinant ratio for this same $f_L$ is
\[
 \frac{\det H_N^{\Phi_L\sigma}}{\det H_N^\sigma}
 =\prod_{a=0}^N\prod_{j=0}^{2L-1}(a+j+1/2).
\]
For clarity its monic norms follow with the original mass as follows.
Gamma recurrence gives
$\Phi_L(y)d\sigma(y)=4^L|\Gamma(L+1/4+iy/2)|^2dy/(2\pi)$.
Write $\alpha=L+1/4$. The Euler integral and its Fourier
autocorrelation give the exact tilted mass
\[
 \int e^{ty}\Phi_L(y)d\sigma(y)
       =\sqrt2\Gamma(2\alpha)(\cos t)^{-2\alpha}
       \qquad(|t|<\pi/2).
\]
Indeed, for $g_\alpha(x)=e^{\alpha x-e^x}$, its Fourier transform is
$\Gamma(\alpha+i\xi)$; the autocorrelation integral at displacement
$u$ is $\Gamma(2\alpha)/(2\cosh(u/2))^{2\alpha}$ after $v=e^x$.
Fourier inversion, $y=2\xi$, and analytic continuation through the
strip justified by the Euler contour bound give the displayed formula,
including its factor $4^L2^{1-2\alpha}=\sqrt2$.
Define the monic polynomials by
\[
 \sum_{j\ge0}p_j^{(\alpha)}(y)\frac{z^j}{j!}
       =(1+z^2)^{-\alpha}e^{y\arctan z}.
\]
Integrating the product of two such series and using
\[\cos(\arctan z+\arctan w)=(1-zw)/\sqrt{(1+z^2)(1+w^2)}\]
gives $\sqrt2\Gamma(2\alpha)(1-zw)^{-2\alpha}$.
For sufficiently small real $z,w$, differentiated integration is
justified by the same exponential bound. Coefficient comparison
proves the squared norm $\sqrt2\,j!\Gamma(j+2L+1/2)$ and all
off-diagonal inner products zero. At $L=0$ the same norm is
$\sqrt2\,j!\Gamma(j+1/2)$. Their ratio is exactly the displayed
finite product. Multiplication by $i^j$ and $y=(S-c)/i$ makes these
polynomials monic in $S$ without changing their norms; all monic
coefficient transformations have determinant one. This proves the
source determinant identity, and cancellation gives precisely the
two source product ranges in MFG5. Thus the actual homogeneous
four-endpoint return for this independent Gamma order is
\[
 \boxed{T_{q,L}^{\rm hom}
       =P_{q,L}-\log r_{q,L}+\log(Z_{q,L}/Z_{q,0}).}
 \tag{MFG6}
\]
This finite identity uses $\Phi_L$ throughout; it does not substitute
the weight $y^{2L}$ for MFG2.

\section{The actual macroscopic potential from the full polynomial}

Write $\lambda_q=L/q$ and suppose $\lambda_q\to\lambda\in(0,1/4]$.
For the literal scale $t=q^2x$ define
\[
 \phi_{q,L}(x)=\prod_{j=0}^{L-1}
             \left(x+\left(\frac{2j+1/2}{q}\right)^2\right),
 \qquad \Phi_L(q\sqrt x)=q^{2L}\phi_{q,L}(x).
 \tag{MFG7}
\]
Define the explicit integral
\[
 h_\lambda(x)=\int_0^\lambda\log(x+4u^2)\,du
 =\lambda\log(x+4\lambda^2)-2\lambda
                +\sqrt x\arctan\frac{2\lambda}{\sqrt x}.
 \tag{MFG8}
\]
The equality follows by differentiating the displayed antiderivative
in its upper limit, with value zero at that limit zero.

For fixed $x>0$, the integrand in MFG8 is increasing with $u\ge0$.
The point $(j+1/4)/q$ lies in the $j$th cell of width $1/q$.
Both the sum of its values and $q$ times the integral lie between
the same left and right sums. Their difference is therefore bounded
by the endpoint difference. This proves the global finite estimate
\[
 \left|\log\phi_{q,L}(x)-q h_{\lambda_q}(x)\right|
 \le\log\left(1+\frac{4\lambda_q^2}{x}\right).
 \tag{MFG9}
\]
In particular the convergence to $h_\lambda$ after division by $q$
is uniform on every compact subinterval of $(0,\infty)$.

The exact macroscopic field is consequently
\[
 \boxed{V_\lambda(x)=\pi\sqrt x-2\log x-2h_\lambda(x).}
 \tag{MFG10}
\]
For $\lambda>0$ it has the equivalent derivative
\[
 V_\lambda'(x)=\frac{\arctan(\sqrt x/(2\lambda))}{\sqrt x}-\frac2x
              =\int_{2\lambda}^{\infty}\frac{du}{x+u^2}-\frac2x.
 \tag{MFG11}
\]
Indeed $\partial_xh_\lambda(x)=
\arctan(2\lambda/\sqrt x)/(2\sqrt x)$, and the two positive
arctangents sum to $\pi/2$. At $\lambda=0$, MFG10 returns to
$V_0(x)=\pi\sqrt x-2\log x$. The extra term in MFG10 has order one
in this macroscopic field; removing it would change the variational
quantity being calculated.

For precision define
\[
 F_\lambda=\sup_{\eta\in\mathcal P((0,\infty))}
 \iint\left[\log|x-y|-\frac{V_\lambda(x)+V_\lambda(y)}2\right]
                         d\eta(x)d\eta(y).
 \tag{MFG12}
\]
The definition is the single bounded-above kernel integral. The explicit
macroscopic equilibrium provider KME1--24 in
\texttt{../equilibrium/MACROSCOPIC\_MARKED\_EQUILIBRIUM.tex}, with
its complete \texttt{endpoint\_existence.tex} KEP1--20, constructs
its unique positive density
$\rho_\lambda$ on $[a_\lambda,b_\lambda]\subset(0,\infty)$, proves
the full exterior variational inequality, and proves finite entropy
relative to the measure $\omega$ below. Its endpoint system is
\[
 \int_{2\lambda}^{\infty}
       \frac{du}{\sqrt{(a_\lambda+u^2)(b_\lambda+u^2)}}
       =\frac2{\sqrt{a_\lambda b_\lambda}},
\]
\[
 \int_{2\lambda}^{\infty}\left[1-
       \frac{u^2}{\sqrt{(a_\lambda+u^2)(b_\lambda+u^2)}}\right]du=4,
\]
\[
 \rho_\lambda(x)=\frac{\sqrt{(b_\lambda-x)(x-a_\lambda)}}{2\pi x}
 \int_{2\lambda}^{\infty}
       \frac{u^2\,du}{(x+u^2)
              \sqrt{(a_\lambda+u^2)(b_\lambda+u^2)}}.
 \tag{MFG13}
\]
The symbol $\rho_\lambda$ here denotes exactly that provider's
$\rho_\kappa$ under $\kappa=2\lambda$; the original centre remains
$c=k/2$. KME15 proves its full variational inequality and KME16--17
prove its compact-support and entropy properties.
The endpoint existence calculation supplies the unique ordered solution,
including $\lambda=0$. At zero this is exactly the original equilibrium
with $uK(\kappa)=2$, $vE(\kappa)=4$. The following argument proves the
new ensemble passage from the literal polynomial to this field.

\section{The full original ensemble, including the discrete factors}

Retain the exact original density map used in GEL3:
\[
 d\omega(x)=\tfrac12x^{-1/2}e^{-\sqrt x}\,dx,
 \quad B_q(x)=2qe^{\sqrt x}\sigma(q\sqrt x),
 \quad d\nu(q^2x)=B_q(x)d\omega(x).
\]
The measure $\omega$ has mass one. Expansion of the two Vandermonde
determinants and the substitution $t_i=q^2x_i$ give
\[
 \boxed{H_s^{(a;L)}
 =q^{2as+2s(s-1)+2Ls}J_{s,q,L}(a),}
\]
\[
 J_{s,q,L}(a)=\frac1{s!}\int
       \prod_{i<j}(x_i-x_j)^2
       \prod_{i=1}^s x_i^a\phi_{q,L}(x_i)B_q(x_i)\,d\omega(x_i).
 \tag{MFG14}
\]
The exact powers come from the original $t^a$ factors, squared
Vandermonde differences, and the $L$ original quadratic factors per
variable, respectively. Every original mass and Jacobian remains in
$B_qd\omega$.

For sequences with $s\to\infty$, $q/s\to2$, $a/s\to2$ and
$L/q\to\lambda\in(0,1/4]$, we prove
\[
 \boxed{\frac{\log J_{s,q,L}(a)}{s^2}\longrightarrow F_\lambda.}
 \tag{MFG15}
\]
We use the previously proved original Gamma inequalities, valid at every
real $y$, with $C_\Gamma=\Gamma(1/4)^2/(2\pi)$:
\[
 C_\Gamma(1+4y^2)^{-3/4}e^{-\pi|y|/2}
 \le\sigma(y)\le
 C_\Gamma e^{1/2}(1+4y^2)^{1/4}e^{-\pi|y|/2}.
 \tag{MFG16}
\]
In particular, on every fixed compact interval $[a,b]\subset(0,\infty)$,
$\log B_q(x)=-\pi q\sqrt x/2+O(\log q)$ uniformly. Globally the
same bounds give
\[
 B_q(x)\le C_q\exp[-(\pi q/2-3/2)\sqrt x],
 \qquad C_q=\frac{\Gamma(1/4)^2}{\pi}
                     q e^{1/2}(1+2q)^{1/2}.
 \tag{MFG17}
\]

\paragraph{Upper bound with the singular endpoints retained.}
Put $R(x)=1+\sqrt x+|\log x|$. MFG9 gives a sequence
$\varepsilon_q\to0$ such that, globally on $x>0$,
\[
 \left|q^{-1}\log\phi_{q,L}(x)-h_\lambda(x)\right|
                  \le\varepsilon_q R(x).
 \tag{MFG18}
\]
To check this global assertion, the endpoint error in MFG9 is at most
$C+|\log x|$ for $0<x\le1$ and at most a constant for $x\ge1$.
For all sufficiently large $q$, $\lambda_q$ lies in a fixed positive
compact interval around $\lambda$. The upper-limit derivative
$\log(x+4u^2)$ in that interval is bounded in absolute value by
$C+\log(1+x)\le C+2\sqrt x$. The difference
$h_{\lambda_q}-h_\lambda$ therefore satisfies MFG18 as well.
Also $|h_\lambda(x)|\le C R(x)$, directly from MFG8.

After removing $C_q^s/s!$ by MFG17, the exponent of the remaining
integral is
\[
 2\sum_{i<j}\log|x_i-x_j|-s\sum_i V_s(x_i),
\]
\[
 V_s(x)=\left(\frac{\pi q}{2s}-\frac3{2s}\right)\sqrt x
        -\frac as\log x-\frac1s\log\phi_{q,L}(x).
\]
The parameter limits and MFG18 show
$|V_s-V_\lambda|\le o(1)R$. For $s>1$, this exponent equals
$\sum_{i\ne j}K_{Q_s}(x_i,x_j)$, where $Q_s=sV_s/(s-1)$ and
$K_Q(x,y)=\log|x-y|-(Q(x)+Q(y))/2$. Therefore for each sufficiently
small fixed $\varepsilon>0$, it is bounded above eventually by
$\sum_{i\ne j}K_{V_\lambda-\varepsilon R}(x_i,x_j)$.

Here is the complete truncation step for this field. For
$0\le\varepsilon\le\varepsilon_0<\min(1,\pi/2)$, put
$Q_\varepsilon=V_\lambda-\varepsilon R$ and
$r_\varepsilon(x)=\log(1+x)-Q_\varepsilon(x)/2$.
The value $h_\lambda(x)$ has a finite limit at zero and grows only
logarithmically at infinity. Thus $r_\varepsilon$ tends to $-\infty$
at both endpoints, uniformly for this interval of $\varepsilon$,
and is uniformly bounded above. The inequality
$K_{Q_\varepsilon}(x,y)\le r_\varepsilon(x)+r_\varepsilon(y)$ proves
that
\[
 K_{\varepsilon,M}(x,y)=\max(K_{Q_\varepsilon}(x,y),-M)
\]
extends continuously to the compact square $[0,\infty]^2$; its value
on the diagonal and wherever a coordinate is an endpoint is $-M$.
Let $\Phi_{\varepsilon,M}$ be its largest integral over probability
measures on that compact interval. For the empirical measure
$\eta_s=s^{-1}\sum_i\delta_{x_i}$,
\[
 \sum_{i\ne j}K_{Q_\varepsilon}(x_i,x_j)
 \le s^2\iint K_{\varepsilon,M}\,d\eta_s^2+sM
 \le s^2\Phi_{\varepsilon,M}+sM.
\]
Because $\omega$ has mass one and
$\log(C_q^s/s!)/s^2\to0$, this proves
$\limsup s^{-2}\log J_{s,q,L}(a)\le\Phi_{\varepsilon,M}$.

For completeness,
$\Phi_{\varepsilon,M}\downarrow F(Q_\varepsilon)$ as $M\to\infty$.
Choose maximizing probabilities for increasing $M$ and pass to a weak
subsequence on the compact interval. For every fixed truncation level,
continuity of its kernel and its domination of subsequent truncations
show that the limiting probability has untruncated kernel integral at
least $\lim_M\Phi_{\varepsilon,M}$. Monotone convergence applies to
a common upper bound minus the kernels. This integral is finite from
below, because a compact interval probability with finite logarithmic
energy is a competitor of finite value at every truncation. It follows
that the limiting probability has no mass at the two endpoints and
gives a legitimate competitor for $F(Q_\varepsilon)$. The reverse
inequality follows by testing every such competitor before truncation.

Finally $F(Q_\varepsilon)\to F_\lambda$ as $\varepsilon\downarrow0$.
The inequality $F(Q_\varepsilon)\ge F_\lambda$ follows from $R\ge0$.
For each fixed $M$, the truncated kernels converge uniformly to
$K_{0,M}$: uniform coercivity puts all pairs with a coordinate outside
a fixed compact interval at value $-M$, and on the remaining compact
square the boundedness of $R$ gives uniform convergence. Thus
$\limsup_{\varepsilon\downarrow0}F(Q_\varepsilon)
\le\Phi_{0,M}$ for every $M$. Letting $M\to\infty$ gives the claim.
These steps prove the required upper bound in MFG15.

\paragraph{Lower bound on the proved actual equilibrium.}
Use the probability $\rho_\lambda(x)dx$ from MFG13. Its support is
compact inside $(0,\infty)$, its logarithmic energy is finite, and its
entropy relative to $\omega$ is finite, as proved by the explicit
equilibrium construction. Restrict the ensemble integral to that
support and apply Jensen's inequality after changing each $d\omega$
to $\rho_\lambda dx$. Writing $f=d(\rho_\lambda dx)/d\omega$, the
result is
\[
 \begin{aligned}
 \log J_{s,q,L}(a)\ge{}&-\log(s!)
 +s(s-1)\iint\log|x-y|\rho_\lambda(x)\rho_\lambda(y)\,dx\,dy\\
 &+as\int\log x\,\rho_\lambda(x)dx
 +s\int\log\phi_{q,L}(x)\rho_\lambda(x)dx\\
 &+s\int\log B_q(x)\rho_\lambda(x)dx
 -s\int\log f(x)\rho_\lambda(x)dx.
 \end{aligned}
\]
All terms are integrable under these proved density properties.
Divide by $s^2$, apply the compact-uniform limits MFG16 and MFG18,
and use $q/s\to2$, $a/s\to2$. The right side tends to
the energy of $\rho_\lambda$ for $V_\lambda$, which equals
$F_\lambda$ by its full variational inequality. This proves the lower
bound and completes MFG15 for all four moving sizes and exponents
needed in MFG4.

\section{Both high endpoints, the source products, and the actual low factor}

Apply MFG14--15 to $(s,a)=(n,q),(n+1,q),(n,q+1),(n,q+1)$.
The same proof covers their $L=0$ counterparts with the moving
exponents $a/s\to2$: set $\phi_{q,0}=1$ and $h_0=0$ identically.
Then the upper field still satisfies $|V_s-V_0|\le o(1)R$, and
the lower Jensen bound uses the already proved $\rho_0$. This gives
the same value $F_0=\mathcal F(2,\pi)$ as GEL10 while explicitly
retaining these moving denominator exponents. Subtract these four
denominator limits. The exact new scale power
is $2L[n+(n+1)+n+n]=2L(2q+1)$. The four squared sizes divided by
$q^2$ sum to a quantity tending to one. Therefore
\[
 \boxed{\log(Z_{q,L}/Z_{q,0})
 =2L(2q+1)\log q+q^2(F_\lambda-F_0)+o(q^2).}
 \tag{MFG19}
\]

Define the explicit source integral
\[
 Q(\lambda)=2\int_1^2\int_0^{2\lambda}\log(s+t)\,dt\,ds
 =2\big[A(2+2\lambda)-A(1+2\lambda)-A(2)+A(1)\big],
\]
\[
 A(v)=\frac{v^2}{2}\log v-\frac{3v^2}{4},\qquad A''(v)=\log v.
 \tag{MFG20}
\]
The closed form follows by integrating the two literal endpoint ranges.
There are exactly $4Lq$ original factors in MFG5. After extracting
$q$ from each factor, use cells of side $1/q$ in
$[1,2]\times[0,2L/q]$. The first range samples the left $s$ endpoint,
the second the right $s$ endpoint, and both sample the midpoint in $t$.
Both partial derivatives of $\log(s+t)$ have absolute value at most
one there. Each individual cell error, after multiplication by $q^2$,
is at most $3/(2q)$. Summing over the $2Lq$ cells in each range gives
the finite bound
\[
 \boxed{\left|P_{q,L}+4Lq\log q+q^2Q(\lambda_q)\right|\le6L.}
 \tag{MFG21}
\]
This preserves the half-integer shifts and both unequal endpoint ranges.

The low factor has an explicit finite control which also preserves
the actual polynomial. The original full-source bounds are
\[
 C_Le^{-\pi|y|}\le\sigma(y)\le C_Ue^{-|y|},\qquad
 C_L=\frac{\pi}{\Gamma(3/4)^2},\quad
 C_U=\frac{\Gamma(1/4)^2}{2\pi\sqrt{\cos1}}.
\]
Define
\[
 A_{q,L}=\int_0^\infty z^{2q}\phi_{q,L}(z^2)\sigma(qz)\,dz,
 \qquad B_q^{\rm low}=\int_0^\infty z^{2q}\sigma(qz)\,dz.
\]
Evenness and the exact substitution $y=qz$ give
$r_{q,L}=q^{2L}A_{q,L}/B_q^{\rm low}$; both half-axis factors and
Jacobians cancel only in this particular ratio.
On $1\le z\le2$, all factors of $\phi_{q,L}(z^2)$ are at least one,
so both $A_{q,L}$ and $B_q^{\rm low}$ are at least $C_Le^{-2\pi q}$.
Since $(2j+1/2)/q\le1/2$ for every stipulated $j<L\le q/4$,
$\phi_{q,L}(z^2)\le(1+z)^{2L}$ on the whole positive axis.
For $D=2q+2L\le5q/2$, the elementary bounds
\[
 (1+z)^D\le2^D(1+z^D),\qquad D!\le D^D
\]
give
\[
 A_{q,L}\le\frac{2C_U}{q}5^{5q/2},\qquad
 B_q^{\rm low}\le\frac{C_U}{q}4^q.
\]
These are obtained by integrating $1$ and $z^D$ against $e^{-qz}$;
their exact integrals are $1/q$ and $D!/q^{D+1}$. Thus
\[
 \boxed{\begin{aligned}
 \log\frac{C_L}{C_U}+\log q-(2\pi+\log4)q
 \le{}&\log r_{q,L}-2L\log q\\
 \le{}&\log\frac{2C_U}{C_L}-\log q
                      +(2\pi+\tfrac52\log5)q.
 \end{aligned}}
 \tag{MFG22}
\]
The remaining low logarithm is therefore $O(q)$ uniformly on the stated
macroscopic order range. No power approximation to $\Phi_L$ was used.

\section{The evaluated macroscopic coefficient on the actual polynomial source}

Substitute MFG19, MFG21 and MFG22 into the exact MFG6. The scale
terms cancel with their full integer coefficients:
\[
 -4Lq\log q+2L(2q+1)\log q-2L\log q=0.
\]
Continuity of the explicit $Q$ and $\lambda_q\to\lambda$ give
\[
 \boxed{\frac{T_{q,L}^{\rm hom}}{q^2}\longrightarrow
 \Psi(\lambda):=F_\lambda-F_0-Q(\lambda),
 \qquad 0<\lambda\le1/4.}
 \tag{MFG23}
\]
In particular $L=q/4$, $K=q+1$ gives the exact coefficient
$\Psi(1/4)$, defined by the proved field MFG10 and the elementary
closed expression MFG20. This is the $q^2$ term of the actual Gamma
polynomial source, rather than the coefficient obtained by changing
that polynomial to a power.

For comparison, the exact identity connecting the two source weights is
\[
 \Phi_L(y)=|y|^{2L}
             \prod_{j=0}^{L-1}\left(1+\frac{(2j+1/2)^2}{y^2}\right)
 \quad(y\ne0).
 \tag{MFG24}
\]
On $y=q\sqrt x$ the logarithm of the second factor divided by $q$
tends to $h_\lambda(x)-\lambda\log x$. Hence the homogeneous power
field $\pi\sqrt x-(2+2\lambda)\log x$ and the actual field MFG10
are related by the exact limiting density correction
$-2[h_\lambda(x)-\lambda\log x]$. The latter correction remains
in MFG23 at fixed positive $\lambda$.

The derivative proved in KME24 by testing the two actual minimizers
in each other's variational problems is
$F'_\lambda=2\int\log(x+4\lambda^2)\rho_\lambda(x)dx$.
Since $Q'(\lambda)=4\int_1^2\log(s+2\lambda)ds$, it gives the exact
integral form
\[
 \Psi(\lambda)=\int_0^\lambda\left[
 2\int\log(x+4t^2)\rho_t(x)dx
                   -4\int_1^2\log(s+2t)ds\right]dt.
 \tag{MFG25}
\]
At zero the endpoint family returns to the original equilibrium.
Consequently $\Psi'(0)=2\mathcal L-4(2\log2-1)=\mathcal C_\Gamma$,
the previously proved small-order coefficient, which is strictly
positive. This establishes the exact local bridge to the earlier
order-$lq$ calculation while retaining the finite positive-$\lambda$
field needed for $K/q\to1$.

\paragraph{Dependencies and the original receiving map.}
The Gamma bounds and the unmodified $L=0$ ensemble limit are GEL1--10.
The constructive actual equilibrium and its full variational identities
are the companion macroscopic equilibrium provider KME1--24, with
complete endpoint existence KEP1--20. Its KME22a evaluates each
$F_\lambda$ by real one-dimensional endpoint integrals, retaining all
additive constants in MFG10. The argument above proves the discrete polynomial
limit, the complete ensemble passage, all scale cancellations, and the
finite low and source controls. Transfer from $\chi_0$ to the original
displaced $\chi$ uses its original source comparison with $f_L$ kept
intact. The baseline at the same independent order $K$, its arithmetic
deficit, and the mixed kernel/boundary maps must then enter their exact
receiving equality at that same $K$.

